\item Está leyendo su libro de texto sobre la mitología griega y encuentra una historia sobre Dédalo e Ícaro. Dédalo construyó dos juegos de alas con plumas y cera, una para él y otra para su hijo Ícaro. El padre y el hijo planearon usar las alas para escapar de su encarcelamiento en la isla de Creta. El padre advirtió a Ícaro que no volara demasiado alto porque la proximidad al Sol podría derretir la cera en sus alas. Por supuesto, Ícaro fue alcanzado por la emoción de volar y voló muy cerca del Sol. Sus alas se derritieron y él cayó al mar. Al leer esta información, piensa en su clase de física, donde su instructor acaba de hablar sobre la temperatura de equilibrio de un objeto sin atmósfera a una distancia dada del Sol. Busca en sus notas y encuentra la siguiente ecuación para esta temperatura de equilibrio:
$$ T = (255\text{ K}) \sqrt{\frac{R}{r}} $$
donde $R$ es la distancia del Sol a la Tierra, $r$ es la distancia del Sol al objeto y $T$ está en Kelvin. Esto plantea un acertijo en su mente: si Ícaro voló tan cerca del Sol que la cera de sus alas se derritió, ¿todavía habría aire en esa ubicación para permitirle volar a esa posición? Tome el punto de fusión de la cera como $65^\circ\text{C}$.
